% Scope: Develop selective excitation and constrained sequence synthesis using the established RF conventions; distinguish mathematical designs from scanner-feasible waveforms.
% Status: architecture only; no chapter prose or completed problems yet.
% Prerequisites: Chapters 4, 5, 8 and 17.
% Planned worked examples: Time-bandwidth product, slice thickness and minimum gradient timing.
% Planned figures: Sinc profiles, excitation k-space and gradient ramps.
% Planned challenge problems: Design a slice-selective pulse; compare SAR and duration under hardware limits.
\chapter{RF Pulses and Sequence Design}
\label{ch:18}

\section{Selective excitation theory}
\label{sec:18:selective-excitation-theory}
\subsection{Small-tip-angle integral and excitation k-space}
\subsection{Sinc pulses, apodization and slice profiles}
\subsection{RF bandwidth and time-bandwidth product}

\section{Large-tip and robust pulses}
\label{sec:18:large-tip-and-robust-pulses}
\subsection{Composite rotations and error compensation}
\subsection{Adiabatic following and inversion}
\subsection{SLR design and excitation-refocusing differences}

\section{Multidimensional and multiband excitation}
\label{sec:18:multidimensional-and-multiband-excitation}
\subsection{Spatially selective trajectories}
\subsection{Multiband phases and peak RF power}
\subsection{Parallel transmit design and local heating constraints}

\section{Joint RF and gradient design}
\label{sec:18:joint-rf-and-gradient-design}
\subsection{VERSE and waveform reparameterization}
\subsection{Off-resonance and relaxation limitations}
\subsection{Gradient amplitude, slew and raster constraints}

\section{Sequence timing and gradient moments}
\label{sec:18:sequence-timing-and-gradient-moments}
\subsection{Trapezoid design and minimum duration}
\subsection{Rephasing, crushers and flow compensation}
\subsection{TE, TR and echo-spacing feasibility}

\section{Design verification}
\label{sec:18:design-verification}
\subsection{Bloch simulation across position and detuning}
\subsection{Slice-profile and pathway validation}
\subsection{SAR implications and scanner implementation checks}

\section{Worked examples}
\label{sec:18:worked-examples}
% Planned calculations: Time-bandwidth product, slice thickness and minimum gradient timing.
\subsection{Derivation and quantitative calculation}
\subsection{Assumptions, limiting cases and scanner interpretation}

\section{Challenge problems}
\label{sec:18:challenge-problems}
% Problem design: Design a slice-selective pulse; compare SAR and duration under hardware limits.
% Use challengeproblem and stable labels prob:18:<topic>.
% Complete solutions belong in Appendix E, section for this chapter.
% Use \begin{solution}{prob:18:<topic>} ... \end{solution} there.
\subsection{Derivation and quantitative design}
\subsection{Model criticism and integrated reasoning}
